\section{Analytical Methodology}
\label{sec:methodology}

\subsection{Data Preprocessing and Feature Engineering}
\label{sec:features}

Preprocessing is deterministic and uses only information that is available
when the forecast for hour \(t\) is issued. Weather at hour \(t\) is treated
as a forecast the operator already holds. Demand features use hours strictly
before \(t\).

\paragraph{Aggregation.}
\texttt{src/data\_loader.py} downloads the UCI file, parses the timestamp, sums
the three zones into \texttt{total\_kw}, and resamples to one-hour bins.
Weather columns and zone loads are averaged inside the hour.

\paragraph{Calendar features.}
From the timestamp the pipeline builds hour of day (\(0\)--\(23\)), day of
week (\(0\)--\(6\)), month, and a weekend flag \texttt{is\_weekend}. These
variables let a tree model represent the daily and weekly cycle without a
separate seasonal decomposition.

\paragraph{Lag features.}
Demand at \(t-1\,\mathrm{h}\), \(t-2\,\mathrm{h}\), \(t-24\,\mathrm{h}\) and
\(t-168\,\mathrm{h}\) is added so the model can use short-run persistence, the
previous day, and the same hour of the previous week. The weekly lag is also
the forecast of the seasonal-naive baseline.

\paragraph{Rolling-window features.}
A trailing mean and a trailing standard deviation are computed on windows of
3, 6 and 24 hours. Each window is applied to the series shifted by one hour.
The value at time \(t\) therefore does not enter its own features. This is
the leakage control described in Section~\ref{sec:leakage}.

\paragraph{Weather and interaction features.}
Temperature, humidity, wind speed and the two diffuse-radiation channels are
kept as exogenous inputs. Because the link between temperature and load is
curved (cooling load rises once the day is hot), the feature set also
includes a cooling degree above \(22^\circ\mathrm{C}\), squared temperature,
and the product of temperature and hour. Tree ensembles can split on these
terms directly. The linear ridge model receives the same columns after
standardisation, so the comparison stays on one feature matrix.

\subsection{Exploratory Data Analysis}
\label{sec:eda}

Exploration is done on the hourly series before any model is scored on the
test block. Four views are required. Each figure below is a placeholder for
the PNG written by \texttt{src/evaluation.py} into \texttt{outputs/figures/}.
Replace the box with \verb|\includegraphics| once the file exists.

\begin{figure}[h]
  \centering
  \figplaceholder{Line plot: mean demand by hour, weekday and weekend\\
  (\texttt{daily\_profile.png}), plus a two-week trace (\texttt{weekly\_series.png})\\
  and monthly means (\texttt{monthly\_mean.png}).}
  \caption{Demand level across the day, a sample of weeks, and the year.}
  \label{fig:seasonality}
\end{figure}

The daily profile shows when the city peaks. The weekly trace shows the
repeat of that shape from one day to the next. The monthly bars show the
seasonal lift, which in this climate is expected to follow temperature.

\begin{figure}[h]
  \centering
  \figplaceholder{Scatter of temperature against demand (\texttt{temp\_scatter.png})\\
  and the weather correlation heatmap (\texttt{correlation\_heatmap.png}).}
  \caption{Association between weather and city demand.}
  \label{fig:weather}
\end{figure}

The scatter is the check for a non-linear temperature response. The heatmap
ranks linear association of temperature, humidity, wind and radiation with
demand. Record the Pearson correlations from \texttt{profile.corr\_with\_demand}
in the text under the figure after the run. Temperature is the weather
variable the operating discussion relies on.

\begin{figure}[h]
  \centering
  \figplaceholder{Autocorrelation bars for lags 0 to 168 hours\\
  (\texttt{autocorrelation.png}).}
  \caption{Autocorrelation of hourly city demand.}
  \label{fig:acf}
\end{figure}

The autocorrelation plot justifies the lag set. Peaks near 24 hours and
168 hours support the daily and weekly features. The values at lags 1, 2, 24
and 168 are stored in \texttt{profile.acf}.

\subsection{Modeling and Validation Approach}
\label{sec:models}

The supervised table is cut into three contiguous blocks in time order. There
is no shuffle and no overlap. With the default ratios the earliest 70\% is
the training block \(D_{\mathrm{train}}\), the next 15\% is the validation
block \(T_{\mathrm{valid}}\), and the last 15\% is the test block
\(T_{\mathrm{test}}\). Hyperparameters are chosen by validation MAE.
\(T_{\mathrm{test}}\) is scored once, after that choice.

\begin{table}[h]
  \centering
  \small
  \begin{tabular}{@{}llll@{}}
    \toprule
    Block & Role & Rows & Span \\
    \midrule
    \(D_{\mathrm{train}}\) & Fit the model & [XX] & [start] to [end] \\
    \(T_{\mathrm{valid}}\) & Select hyperparameters & [XX] & [start] to [end] \\
    \(T_{\mathrm{test}}\) & Report MAE, RMSE, WAPE, peak scores & [XX] & [start] to [end] \\
    \bottomrule
  \end{tabular}
  \caption{Chronological backtest. Fill the cells from \texttt{metrics.json} key \texttt{splits}.}
  \label{tab:splits}
\end{table}

Five forecasters share that split.

\begin{enumerate}[label=(\arabic*), itemsep=2pt]
  \item \textbf{Seasonal naive.} The forecast for hour \(t\) is the observed
        demand at \(t-168\). This baseline is strong whenever the weekly cycle
        dominates, and it is the reference for every learned model.
  \item \textbf{Ridge regression.} A standardised linear model. The penalty
        \(\alpha\) is chosen from \(\{0.1, 1, 10, 100\}\) on
        \(T_{\mathrm{valid}}\).
  \item \textbf{Random forest} \cite{breiman2001}. Two depths are compared on
        the validation block. The forest captures threshold effects of
        temperature and hour that a single linear coefficient cannot.
  \item \textbf{XGBoost} \cite{chen2016}. Gradient boosting with early
        stopping on \(T_{\mathrm{valid}}\). Learning rate and tree depth are
        selected by validation MAE.
  \item \textbf{LightGBM} \cite{ke2017}. The same protocol as XGBoost, with
        the number of leaves in place of depth. LightGBM is the model used
        for the feature-importance discussion and for the default dashboard
        view.
\end{enumerate}

The selected parameter dictionary for each model is written to
\texttt{selected\_params} in the metrics file. Paste that dictionary into
the appendix or under Table~\ref{tab:scores} when the report is finalised.
